% Additive human-facing mathematical route; fixed accepted 692-page baseline.
\section{The question and the present answer}\label{the-question-and-the-present-answer}

The programme asks whether a cohomological construction over a proposed absolute base can retain the arithmetic information of the Riemann zeta function and provide sufficiently strong control of a specified spectral action. The current calculations study the consequences of a hypothetical off-critical-line zero quartet. They build the associated spaces and maps, preserve repeated zeros, and calculate norms and determinant returns on those actual spaces. This is research toward a geometric argument about zeta zeros. The edition proves a substantial collection of finite-dimensional identities and estimates. It does not prove the Riemann hypothesis, establish the existence of an off-line quartet, or exclude such a quartet.

The inspiration from Pierre Deligne is precise at the level of ambition: a cohomological realization should come with a structure that controls spectral size. Deligne's theorem for smooth projective varieties over finite fields has an actual variety, an actual Frobenius action and a purity statement. Here the base, support structure, analytic classes, periods and metric comparisons have to be constructed and connected. Their formal resemblance to finite-field geometry is a reason to calculate those connections, not a proof of a purity theorem. The original audit constructs specific scaling and nilpotent maps and records the differences between them. Deligne's \emph{La conjecture de Weil. I}, theorem 1.6, supplies that finite-field statement \cite{Deligne1974WeilI}. His later theory of weights is a distinct development \cite{Deligne1980WeilII}. Neither statement is silently imported into the present analytic construction.

The central allocation problem, retained in the 726-page successor, is concrete. A known total determinant growth splits into two correlated contributions. Their sum has been evaluated. Their separate leading allocation has not. The surviving expression is a signed spectral profile of eight specified four-by-four positive matrices, together with four specified complementary quotient returns. The matrices are computed from the original period, Gamma recurrence, conductor division and lower-evaluation minimum; they are not arbitrary test matrices. A later part of this exposition writes down their exact formula.

There is a second, related calculation concerning the original action's canonical allowance. Strong lower bounds for that allowance have now been proved, including growing polynomial-degree windows. An improvement to the mixed allocation does not automatically improve that action bound: an exact cancellation must be applied first. Keeping these two questions connected through that identity is essential to understanding both what has been achieved and what remains.

\section{What the support base is recording}\label{what-the-support-base-is-recording}

The base notation distinguishes \emph{absence} from a \emph{present coefficient equal to zero}. This project construction is documented in \cite{GlobalizationNote2026,SplitZeroSecondary2026}; its human mathematical background and exact maps are developed in Appendix~\ref{app:scalar-and-base}. For a ring \(R\), the underlying carrier is

\[
G(R)=\{\tau\}\sqcup R^\bullet,\qquad e_R=0_R^\bullet.
\]

The two observations are

\[
(p_R,b_R)(\tau)=(0,0),\qquad
(p_R,b_R)(r^\bullet)=(r,1).
\]

Their joint map is injective. Indeed two present points with the same first coordinate have the same coefficient; an absent and a present point have different second coordinates. Thus \(p_R(e_R)=p_R(\tau)=0\) does not identify \(e_R\) with \(\tau\). The support bit records a distinction that the scalar observation forgets. The proposed absolute base is denoted \(\mathbf F_{1,\tau}=\{\tau,1\}\), with its specified pointed/preaddition structure. A structural map sends its absent element to \(\tau\) and its unit to \(1_R^\bullet\). The structural base map belongs to the pointed/preaddition category; the coefficient adjunction belongs to the diagram category. Their complete definitions and proofs are retained in the original base-formation source.

Why does this distinction survive so far into an analytic calculation? Because a relation may map to zero \emph{at an existing support label}. For a coefficient map \(L\), its lift on a present label \(\omega\) is

\[
(\omega,x)\longmapsto(\omega,Lx).
\]

If \(Lx=0\), the image is \((\omega,0)\), not external absence. When a larger coefficient face is obtained by zero insertion \(j_A^B\), the actual compatibility required is \(L^B j_A^B=j_A^B L^A\). This is a commutative-square statement about the maps. It preserves the original outer leg and every coefficient face; it does not promote a coefficient-linear section to a unital algebra morphism. {[}R09, §1, ``Scope, original input, and retained support''; R14, AKP1--5, §26, pp.110--111. The full categorical definitions and proofs are in the original base-formation audit, chapter 01, §§1--4.{]}

\section{The analytic source: theta, Mellin transform and full jets}\label{the-analytic-source-theta-mellin-transform-and-full-jets}

Write the Riemann xi function in the convention

\[
\xi(s)=\frac12s(s-1)\pi^{-s/2}\Gamma(s/2)\zeta(s),
\qquad g(s)=2\xi(s).
\]

The Riemann zeta and theta setting is classical \cite{Riemann1859}. The original test space and theta map are

\[
V=\{\phi\in\mathcal S(\mathbb R)_{\mathrm{ev}}:
          \phi(0)=0,\ \int_{\mathbb R}\phi(x)\,dx=0\},
\qquad \Theta\phi(x)=2\sum_{n\ge1}\phi(nx).
\]

The target \(\mathscr B\) consists of smooth functions on \(\mathbb R_{>0}\) for which

\[
\sup_{x>0}x^b|D^jF(x)|<\infty
\quad(b\in\mathbb Z,\ j\ge0),\qquad D=-x\partial_x.
\]

Its Mellin transform is \(\mathcal MF(s)=\int_0^\infty F(x)x^s\,dx/x\). Integration by parts, whose two boundary terms vanish by the defining decay at zero and infinity, gives

\[
\mathcal M(DF)(s)=s\mathcal MF(s).
\]

This is why a polynomial in the Euler operator becomes the same polynomial in the arithmetic variable \(s\). The quotient \(Q=\mathscr B/\Theta V\) records analytic classes. The original two-term coefficient resolution is

\[
[V\oplus V\xrightarrow{\ (\phi,\psi)\mapsto
          \Theta(\phi-\widehat\psi)\ }\mathscr B].
\]

Its chart and tensor-resolution identifications are proved in the original audit; the current arithmetic map uses its top tensor cohomology \(Q^{\otimes k}\). {[}R09 §1; R14 AKP1--2.{]}

This is a direct debt to Ralf Meyer's spectral construction, not merely a generic similarity to the literature. In Meyer, \emph{A spectral interpretation for the zeros of the Riemann zeta function}, §3, Definition 3.1 and equation \(11\), \(Zf=\sum_{n\ge1}f(nx)\), the even Schwartz subspace has \(f(0)=\widehat f(0)=0\), and the Poisson identity supplies its two asymptotic conditions. Here \(V\) is that subspace under the stated Fourier convention and \(\Theta=2Z\). In fact the original quotient \(Q\) is topologically and equivariantly Meyer's quotient \(\mathcal H_-^0\), as the explicit comparison below shows. It should not be presented as a newly invented quotient. The source conventions are related as follows: Meyer's completed zeta denoted \(\xi\) does not contain the same \(s(s-1)\) factor as the programme's \(g\), and the temporary generator \(x\partial_x\) in his Proposition 3.2 has the opposite sign from \(D\). His closed-range result is Theorem 3.3. Meyer in turn explicitly discusses Alain Connes's earlier spectral interpretation. These are specific debts to Meyer and the antecedent he identifies in Connes \cite{Meyer2005,Connes1999}. The version used for the convention comparison is Meyer's pinned arXiv version 3, §§2--4.

For completeness, put \(f(t)=F(e^t)\). The source seminorms become \(p_{b,j}(F)=\sup_te^{bt}|f^{(j)}(t)|\). Meyer requires \(e^{st}f(t)\) to be Schwartz for each real \(s\). If \(b_-<s<b_+\) are integers, every polynomial weight \((1+|t|)^n\) is bounded by a constant times \(e^{(b_+-s)t}\) on \(t\ge0\), and by a constant times \(e^{(b_--s)t}\) on \(t\le0\). Leibniz's rule therefore gives

\[
\sup_t(1+|t|)^n|\partial_t^j(e^{st}f)|
\le C\sum_{r=0}^j\binom jr|s|^{j-r}
                  (p_{b_-,r}(F)+p_{b_+,r}(F)).
\]

Conversely \(e^{bt}f^{(j)}=(\partial_t-b)^j(e^{bt}f)\), so each \(p_{b,j}\) is bounded by a finite sum of Meyer's Schwartz seminorms. The identity is consequently a homeomorphism \(\mathscr B\simeq\mathcal H_-\). Since \(2V=V\), its quotient map is the literal identity \(Q\to\mathcal H_-^0,[F]\mapsto[F]\). Both spaces carry \(U_aF(x)=F(x/a)\), so this map also intertwines their dilation representations. This exact precedent determines the credit at the first definition of \(Q\).

Now stipulate a complete zero quartet

\[
\rho=\tfrac12+\delta+i\gamma,\qquad
0<\delta<\tfrac12,\quad\gamma>2,\quad m\ge1,
\] \[
\mathcal R=\{\rho,\bar\rho,1-\rho,1-\bar\rho\},\qquad
h(s)=\prod_{\alpha\in\mathcal R}(s-\alpha)^m.
\]

The hypothesis specifies that \(m\) is the complete order at each member of this packet. Consequently \(v_h=g/h\) is entire at those four points and \(v_h(\alpha)\ne0\). The primary algebra \(E_h=\mathbb C[s]/(h)\) has dimension \(4m\). It retains derivatives through order \(m-1\), not just the four values. Multiplication by

\[
\upsilon_h=j_hv_h\in E_h
\]

is an automorphism because its constant value in each local primary factor is nonzero. Its matrix in raw derivative coordinates is triangular, with entries

\[
\binom rs v_h^{(r-s)}(\alpha)\qquad(r\ge s).
\]

The off-diagonal entries matter whenever \(m>1\). Replacing this matrix by four scalar values would change the map. The original Euler-section proof constructs \(F_h\) with \(\mathcal MF_h=v_h\) and an exact jet section \(\sigma_h:E_h\hookrightarrow Q\). It proves those maps on their stated domains. {[}R09 §1; R14 OCS1, OCS12--18, AKP2; full analytic proof in original audit chapter 04.{]}

\section{Why the sum algebra is smaller than the full tensor product}\label{why-the-sum-algebra-is-smaller-than-the-full-tensor-product}

Fix \(k=4l+1\), \(l\ge1\), and retain

\[
c=k/2,\qquad e_k=1+k(m-1),\qquad
\lambda_{ab}=c+(2a-k)\delta+i(2b-k)\gamma,
\] \[
\chi(S)=\prod_{a,b=0}^{k}(S-\lambda_{ab})^{e_k},\qquad
q=e_k(k+1)^2,\qquad E=\mathbb C[S]/(\chi).
\]

The full ordered tensor product is \(E_h^{\otimes k}\), of dimension \((4m)^k\). The programme does not identify it with the \(q\)-dimensional algebra \(E\). It gives the exact injection

\[
\eta:E\longrightarrow E_h^{\otimes k},\qquad
\eta[P]=\upsilon_h^{\otimes k}P(s_1+\cdots+s_k)1.
\tag{A}
\]

Here is the reason for the exponent \(e_k\), including its constant. In a fixed ordered primary tensor block write \(s_i=\alpha_i+n_i\), where \(n_i^m=0\). The nilpotents commute. Every monomial in \((\sum_i n_i)^{k(m-1)+1}\) has some exponent at least \(m\), so that power vanishes. In the preceding power the coefficient of \(n_1^{m-1}\cdots n_k^{m-1}\) is

\[
\frac{[k(m-1)]!}{[(m-1)!]^k},
\]

which is nonzero over \(\mathbb C\); this is the only surviving monomial of that total degree. The nilpotence index is therefore exactly \(e_k\). The possible sums of the original roots are precisely the distinct \(\lambda_{ab}\). Their real and imaginary parts separately distinguish \(a,b\), since \(\delta,\gamma>0\). The annihilator of the cyclic unit for multiplication by \(\sum_i s_i\) is thus exactly \(\chi\). Polynomial substitution has kernel \((\chi)\), and multiplication by the invertible full unit \(\upsilon_h^{\otimes k}\) preserves injectivity. This proves (A) and the intertwining identity

\[
\eta M_S=\Bigl(\sum_iM_{s_i}\Bigr)\eta.
\]

This calculation illustrates the intended role of an exact map: the smaller cyclic module is a specified submodule of the complete tensor object. It does not erase the other tensor directions. {[}R14 OCS1--3, §1 pp.20--21; CAI10, §198 p.624.{]}

At the function level the corresponding source is

\[
\mathcal V_{h,k}(P)=P(D_1+\cdots+D_k)F_h^{\otimes k},\qquad
q^{(k)}\mathcal V_{h,k}(P)=\sigma_h^{\otimes k}\eta[P].
\]

The relation \(P\mapsto P+\chi\kappa\) changes the representative by an explicit coboundary. Ordered monic division gives polynomials \(Q_a(\mathbf s)\) with

\[
\chi(s_1+\cdots+s_k)=\sum_{a=1}^kh(s_a)Q_a(\mathbf s).
\]

The final remainder vanishes by the annihilator calculation. With the original \(\phi_0\) satisfying \(\Theta\phi_0=h(D)F_h\), placed in the degree-zero plus leg as \((\phi_0,0)\), define

\[
\Psi_\kappa=\sum_{a=1}^k(-1)^{a-1}Q_a(\mathbf D)
\kappa\!\left(\sum_bD_b\right)
\bigl(F_h^{\otimes(a-1)}\otimes\phi_0
                       \otimes F_h^{\otimes(k-a)}\bigr).
\]

In the \(a\)-th summand the tensor differential contributes another sign \((-1)^{a-1}\); the two signs multiply to \(1\). Applying the differential therefore gives exactly \(d\Psi_\kappa=\mathcal V_{h,k}(\chi\kappa)\). This proves representative independence with its actual homotopy. By contrast, lying in a subsequent observation kernel does not make a class a coboundary: the original injection may still carry it to a nonzero arithmetic class. {[}R14 AKP2--5, pp.110--111.{]}

\section{What the period observation sees}\label{what-the-period-observation-sees}

The additional observation uses the original polynomial primitive \(\Phi_h'=h\), \(\Phi_h(0)=0\), a nonzero parameter \(u\), a fixed branch, and the oriented contours

\[
\Gamma_j=-\ell_0+\ell_j,\qquad
\arg\ell_j=\frac{\arg u+\pi+2\pi j}{4m+1}.
\]

On a degree-less-than-\(4m\) representative,

\[
(\Pi[P])_j=\int_{\Gamma_j}e^{\Phi_h(z)/u}P(z)\,dz.
\]

The preceding full period theorem supplies convergence and invertibility on its stated domain. In file 03, PDE1--13, the calculation uses the actual determinant differential equation and its Gamma/Fourier base point; this proof is given directly, rather than invoking an unspecified abstract period theorem. It is not permissible to choose period entries independently: each is this integral of the specified representative. Let \(n=4m\), \(N=n+1\), \(U=M_{\upsilon_h}\) and \(\iota[P]=[P(s_1+\cdots+s_k)]\). The observed map is

\[
L=P_k\Pi^{\otimes k}U^{\otimes k}\iota=\jmath\Lambda,
\quad P_k=\frac1N\sum_{a=0}^{N-1}(C^{-a})^{\otimes k},
\quad K=\ker\Lambda,
\]

where \(\Lambda:E\twoheadrightarrow B\) is onto its actual image and \(\jmath\) includes that image into the ordered period tensor space. The literal cyclic matrix is

\[
(Cx)_j=x_{j+1}-x_1\ (j<N-1),\qquad
(Cx)_{N-1}=-x_1.
\]

With \(\zeta=e^{2\pi i/N}\), \(f_t=(\zeta^{jt}-1)_{j=1}^{N-1}\), one computes \(Cf_t=\zeta^tf_t\). The finite average selects the symmetric tensor words whose sum of indices is zero modulo \(N\). However, these Fourier columns have Gram \(N(I+\mathbf1\mathbf1^*)\), so this coordinate diagonalization is not an assertion that the projector is orthogonal in the original metric. The current source retains the induced full tensor Gram and all its cross terms. {[}R14 OCS2--11, §1 pp.20--23.{]}

For the simple-quartet case \(m=1\), the corresponding period curve is a quadric with an actual cyclic orbit. Its five-member orbit gives

\[
\operatorname{rank}\Lambda=k^2-6k+17,\qquad
\dim K=8k-16.
\]

The actual period calculation supplies a radius \(R_*\) forcing the relevant orbit on its stated branch domain; the metric applications retain \(|u|\ge R_*\), and later pole presentations have their further specified exclusions or sufficient radius. These numbers are not asserted for arbitrary multiplicity. The all-\(m\) geometry retains its own Segre--Veronese maps, nilpotents and stabilizer. {[}R14 OQX3, §3 pp.32--33; PQO/RPC, §§4--5 pp.34--42; OMG, §13 pp.84ff.{]}

Ranks alone do not determine the desired volumes. The latter depend on the period coefficients and on how their kernel sits inside the original source metric. This is the reason the calculation continues after the exact rank theorem.

\section{The arithmetic norm and the two Gamma references}\label{the-arithmetic-norm-and-the-two-gamma-references}

The arithmetic density is

\[
w_h(y)=\frac{|v_h(\tfrac12+iy)|^2}{2\pi},\qquad
m_{h,k}=w_h^{*k}.
\]

For \(P\in\mathbb C[S]\), the finite source norm is

\[
\|P\|_{\mathrm{ar},k}^2
=\int_{\mathbb R}|P(c+iy)|^2m_{h,k}(y)\,dy.
\]

It comes from the same Mellin/theta class; no replacement probability measure is implicit. In particular the full mass is retained. The analytically tractable reference measures are

\[
\sigma(y)=\frac{|\Gamma(\tfrac14+iy/2)|^2}{2\pi},
\quad c_\sigma=\sqrt{2\pi},\quad
c_a=2^{1-2a}\Gamma(2a),
\] \[
r_k^0(y)=\frac{c_\sigma^k}{c_{k/4}}
          \frac{|\Gamma(k/4+iy/2)|^2}{2\pi}.
\]

Their masses are \(c_\sigma\) and \(c_\sigma^k\), respectively. The monic Gamma-polynomial norm convention is \(h_{n,s}=(2\pi)^{s/2}n!(s/2)_n\), with \(s=1\) or \(k\). These are classical Gamma/Meixner--Pollaczek structures, used here with the displayed scaling and mass convention \cite{DLMF18Current,KoekoekLeskySwarttouw2010}.

An exact identity explains why two source orders can be compared on the same arithmetic quotient. Set

\[
b_j=2j+\tfrac12,\qquad
f_l(S)=\prod_{j=0}^{l-1}(S-(c+b_j)),\qquad
C_k=\frac{c_\sigma^k}{c_{k/4}}4^{-l}.
\]

The classical Gamma recurrence \cite{DLMF5Current} gives

\[
\frac{\Gamma(l+1/4+iy/2)}{\Gamma(1/4+iy/2)}
=(i/2)^l\prod_{j=0}^{l-1}(y-ib_j).
\]

Taking modulus on the real line gives the complete equality

\[
r_k^0(y)=C_k|f_l(c+iy)|^2\sigma(y).
\tag{B}
\]

The holomorphic recurrence uses the displayed minus signs before taking the modulus. For example \(k=5\) gives \(f_1(S)=S-3\), \(C_5=16\pi^2/3\). Since \(k\) is odd, each original \(\lambda_{ab}\) has nonzero imaginary part, whereas all roots of \(f_l\) are real. Thus \(f_l\) and \(\chi\) are coprime, multiplication \(M_f:E\to E\) is invertible, and

\[
\eta M_f=M_{f(\sum_i s_i)}\eta.
\]

This supplies the exact arithmetic map behind \(B\). Multiplication by \(f_l\) changes the source representative in a known way; it is not silently identified with the identity. {[}R09 §2, ``The exact two-reference multiplication,'' and §3, ``Joint full jets and the mixed Schur complement.''{]}

Comparison estimates between these norms and the arithmetic source are proved on \emph{all polynomial directions} before taking restrictions or quotients. If \(aH_1\preceq H_2\preceq AH_1\) on one coefficient space, then minimizing each inequality over the same affine fibre proves the corresponding quotient inequality. A bound on selected moments or selected monic test vectors alone would not prove that whole-form comparison. The current arithmetic comparison is AMT, and later source-order refinements retain those same fibres and maps. {[}R14 §8 pp.52ff.; §§192,197,204--211 pp.571ff.,602ff.,648--659.{]}

\section{From a source norm to four determinant returns}\label{from-a-source-norm-to-four-determinant-returns}

For \(N\ge q-1\), the remainder map \(J_N:\mathcal P_N\twoheadrightarrow E\) has kernel \(\chi\mathcal P_{N-q}\). If \(H_N\) is the source Gram in a fixed coefficient frame, its attained quotient Gram is

\[
G_N=(J_NH_N^{-1}J_N^*)^{-1}.
\tag{C}
\]

Indeed \(S_N^{\mathrm{rem}}=H_N^{-1}J_N^*G_N\) is a right inverse of \(J_N\). For \(z\in\ker J_N\), \(z^*H_NS_N^{\mathrm{rem}}=0\). Every lift of a value is the displayed minimum lift plus such a \(z\), so its squared norm is the minimum squared norm plus \(z^*H_Nz\). This proves (C), not merely its positivity. In a monic orthogonal-polynomial frame, with \(b_n=[p_n]_\chi\) and full squared norm \(\omega_n\), the inverse is

\[
G_N^{-1}=\sum_{n=0}^{N}\frac{b_nb_n^*}{\omega_n}.
\]

The first \(q\) remainder columns form a triangular invertible matrix; hence the inverse exists. {[}R14 CAI2, §198 p.622.{]}

Now apply the \emph{actual} observation \(\Lambda:E\twoheadrightarrow B\) and a fixed inclusion \(I:K\hookrightarrow E\). Its kernel metric and attained image metric are

\[
H_{K,N}=I^*G_NI,\qquad
Q_{B,N}=(\Lambda G_N^{-1}\Lambda^*)^{-1},
\] \[
S_N=G_N^{-1}\Lambda^*Q_{B,N},\quad
\Lambda S_N=I_B,\quad I^*G_NS_N=0.
\]

Choose one fixed section \(S_*\) of \(\Lambda\). The frames \([I,S_*]\) and \([I,S_N]\) differ by a triangular matrix with determinant one, since \(S_N-S_*\) has image in \(K\). Orthogonality therefore proves

\[
\det H_{K,N}\det Q_{B,N}
=|\det[I,S_*]|^2\det G_N.
\tag{D}
\]

The frame factor is present; it is not set equal to one. Use the original four endpoints

\[
N_0=q-1,\quad N_1=q,\quad N_2=2q-1,\quad N_3=2q,
\qquad(\varepsilon_0,\varepsilon_1,\varepsilon_2,\varepsilon_3)
=(1,1,-1,-1).
\]

Define the total, kernel and observed-image returns by

\[
\mathcal B_k=\sum_{i=0}^3\varepsilon_i\log\det G_{N_i},\quad
F_K=\sum_i\varepsilon_i\log\det H_{K,N_i},\quad
F_B=\sum_i\varepsilon_i\log\det Q_{B,N_i}.
\]

Taking this exact signed sum in (D) cancels the fixed factor because \(\sum_i\varepsilon_i=0\), and proves \(F_K+F_B=\mathcal B_k\). The signs are part of the definition; they are what makes constant frame changes cancel. This same calculation works for the arithmetic and each Gamma source separately. {[}R14 CAI3--5, pp.622--623; AKP6--10, pp.111--112.{]}

Schur complements, minimum sections, determinant lemmas and their Woodbury form are classical finite-dimensional linear algebra. The contribution to this programme is the calculation of the original coefficient maps, frames and constants to which those identities are applied, and the proof that they give the original arithmetic/cohomological quantities. The classical linear-algebra identities are credited independently of this programme-specific application \cite{BoydElGhaouiFeronBalakrishnan1994,Hager1989Inverse,ShermanMorrison1950Inverse,Woodbury1950Modified}.

\section{The later native metric and its two correlated row energies}\label{the-later-native-metric-and-its-two-correlated-row-energies}

The conductor calculations recast the relevant dual metric in terms of a numerator divided by the actual conductor symbol \(E_A\). Its order \(v=\operatorname{ord}_0E_A\), leading moment \(\mu_v\ne0\) and scalar \(\alpha_A=v!/\mu_v\) remain in every division matrix. The metric is the squared norm of that numerator after minimizing over all allowed lower evaluations. Thus its finite raw data consist of a lower-evaluation matrix \(V\) and a numerator matrix \(T\), both in the original orthonormal Gamma frame.

At one cutoff set \(A=V^*V\), \(K_0=A^{-1}V^*T\). When the next rows \(v_+,t\) are appended, the residual row and its conditioning denominator are

\[
r=t-v_+K_0,\qquad \lambda=1+v_+A^{-1}v_+^*,\qquad
h=\lambda^{-1/2}r.
\]

If the old attained numerator Gram is \(D\), its new value is exactly

\[
D^+=D+h^*h.
\tag{E}
\]

To check \(E\), subtract the old attained minimum. A change \(\eta\) in the old lower coefficient costs

\[
\eta^*A\eta+\|rx-v_+\eta\|^2.
\]

The normal equations give \(\eta=(A+v_+^*v_+)^{-1}v_+^*rx\). Substitution gives \(x^*r^*r x/(1+v_+A^{-1}v_+^*)\). This proves the update with the complete lower conditioning, rather than minimizing the new row independently of the old rows. {[}R14 HJR14--15, §73 pp.304--309.{]}

Let \(B\) now denote the actual onto coefficient map in the native quotient and \(J\) its kernel inclusion. These local symbols belong to HJR and should not be confused with the earlier period-image space \(B\). Set

\[
D_J=J^*DJ,\qquad D_B=(BD^{-1}B^*)^{-1},\qquad
L=D^{-1}B^*D_B.
\]

The same attained-section proof gives the exact orthogonal inverse decomposition

\[
D^{-1}=J D_J^{-1}J^*+L D_B^{-1}L^*.
\]

Consequently define two nonnegative \emph{actual} row energies

\[
x=hJD_J^{-1}J^*h^*,\qquad
y=hLD_B^{-1}L^*h^*.
\]

The determinant lemma and one further quotient minimum give

\[
\log\frac{\det D^+}{\det D}=\log(1+x+y),
\] \[
\log\frac{\det D_J^+}{\det D_J}=\log(1+x),\qquad
\log\frac{\det D_B^+}{\det D_B}
=\log\left(1+\frac{y}{1+x}\right).
\tag{F}
\]

For the last identity, a quotient lift of a value \(a\) is \(La+Jz\) in the stated original coefficient spaces. Its old squared norm is \(a^*D_Ba+z^*D_Jz\), and its additional squared norm is \(|hLa+hJz|^2\). Put \(t=D_J^{1/2}z\) and \(w=hJD_J^{-1/2}\), so \(ww^*=x\). The same one-row normal equation as in (E) gives the minimum additional form \(|hLa|^2/(1+x)\). Its determinant relative to \(D_B\) is \(1+y/(1+x)\), proving the last identity with that denominator. Thus

\[
\log(1+x+y)=\log(1+x)+\log\left(1+\frac{y}{1+x}\right)
\]

is the determinant splitting of this precise metric, not a free numerical partition of a known total. {[}R14 HJR16--18, pp.304--309.{]}

Summing the rows from \(M=q-1\) through (2q-1) with weight one at those two ends and weight two at every interior row gives

\[
S_K=\sum_Mw_M\log\left(1+\frac{y_M}{1+x_M}\right),\qquad
S_R=\sum_Mw_M\log(1+x_M),\qquad
T=S_K+S_R.
\tag{G}
\]

The original dual/coefficient maps connect \(S_K\) to the conductor-transported kernel return; the finite discrepancy from \(F_K\) is explicitly controlled. The complete two-window row total satisfies, on the original simple-quartet fixed-period domain,

\[
\frac{T}{kq}\longrightarrow16C_\partial,
\qquad21.66851180520<16C_\partial<21.66851180521.
\]

The coefficient is the evaluated logarithmic-potential expression

\[
C_\partial=4\log(4/\pi)-10-2F(2,\pi)+4\mathcal L(2,\pi),
\]

where \(F\) and \(\mathcal L\) are the energy and logarithmic moment of the proved equilibrium for \(V_{\alpha,\beta}(x)=\beta\sqrt x-\alpha\log x\). The exact equilibrium and outward certificate are part of the source. For orientation to the classical logarithmic-potential setting, see \cite{SaffTotik1997}; the specific potential, original arithmetic comparison and receiving constants are calculated in R14. This number is the total coefficient. It does not select the separate leading value of either summand in (G). {[}R14 HJR19--25 and subsequent receiving refinements; ECL15--18, §59 pp.243--244; current continuation §``Exact working quantities''.{]}

\section{The precise matrix expression that remains}\label{the-precise-matrix-expression-that-remains}

The latest mixed calculation uses the target degree \(j=k+4\), the original source pairs \((1,1)\) and \((k,j)\), and eight fixed coefficient subspaces in the four-copy target. In the notation of the current prompt, put

\[
I=\operatorname{im}(B_{\rm cof}J_0),\quad
B=\operatorname{im}B_{\rm cof},\quad
S=\operatorname{im}[B_{\rm cof},F_e],\quad J=\ker\mathbb B,
\] \[
K_i=I\cap J,\quad W_i=I+J,\quad W=B+J,\quad
K_c=S\cap W,\quad T_e=S+J.
\]

Here \(F_e\), \(e=0,1\), is the original connecting-map column family at the two low cutoffs. The full connecting constructions and exact coefficient maps are in R14 §§151--159 and §§184--193. The resulting inclusions are the two flags

\[
J\subset W_i\subset W\subset T_e,\qquad
K_i\subset I\subset K_c\subset S.
\]

At each original row let \(D_0\) be its preceding positive numerator form, \(\rho_n\) its lower-evaluation-corrected row and \(\zeta_n\) its complete conditioning scalar. With

\[
\mathcal D_0=\operatorname{diag}(D_0,D_0,D_0,D_0),\qquad
\mathcal U_n=\operatorname{diag}(a_n,a_n,a_n,a_n),
\qquad a_n=\frac{\rho_n}{\sqrt{1+\zeta_n}},
\]

the four-by-four compression associated with a fixed independent frame \(X\) for one of the eight subspaces is

\[
K_{X,n}=\mathcal U_nX(X^*\mathcal D_0X)^{-1}X^*\mathcal U_n^*.
\tag{H}
\]

The empty subspace gives the zero matrix. Formula (H) is unchanged by an invertible change of columns in \(X\). If \(Y=\mathcal U_n\mathcal D_0^{-1/2}\), it is \(Y\Pi_XY^*\), where \(\Pi_X\) is the orthogonal projector onto \(\mathcal D_0^{1/2}\operatorname{im}X\). Hence inclusion of subspaces gives positive matrix order. Also

\[
YY^*=r_nI_4,\qquad
r_n=\frac{\rho_nD_0^{-1}\rho_n^*}{1+\zeta_n},
\qquad0\preceq K_{X,n}\preceq r_nI_4.
\]

The signed threshold profile \(\mathcal P_{k,e}\) takes the four eigenvalues of each matrix, including zeros, applies \(\log\max(1,\lambda)\), and sums over the original rows with their endpoint/interior weights. The positive subspaces are \((T_e,K_c,W_i,K_i)\); the negative ones are \((W,S,J,I)\). The accepted theorem gives

\[
|\mathcal L_e-\mathcal P_{k,e}|
\le\sum_n\vartheta_n\varepsilon(r_n)\le16q_j\log2,
\] \[
\varepsilon(r)=
\begin{cases}
4\log(1+r),&0\le r\le1,\\
4\log\bigl(4r/(1+r)\bigr),&r\ge1.
\end{cases}
\tag{I}
\]

This error is \(o(kq_k)\). It has already been proved for the unchanged original profile. Its proof uses the two nested flags and eigenvalue monotonicity; it assumes no simultaneous diagonalization. {[}R14 NTE1--7, §199 pp.631--632.{]}

The map back to the original allocation is equally important:

\[
4S_{K,j}=\mathcal L_e+V_b+V_{\rm rem}+V_s+V_A.
\tag{J}
\]

The terms are attained quotient returns on the original coefficient spaces: \(V_{\rm rem}\) is the remaining target quotient by \(T_e\); \(V_s\) is the connecting source form on \(S/K_c\); \(V_A\) is the invariant source form on \(I/K_i\) in the target-degree metric; and \(V_b\) is the retained boundary return. Each is computed with the original four endpoints and fixed value frame. These positive returns are correlated with \(\mathcal L_e\), since all use the same original row data. The degree-\(k\) proper-source form has its own centre and comparison map; it is not substituted for \(V_A\). {[}R14 CPQ1--20, §194 pp.587--591; EFP, §§204--211.{]}

Thus the remaining target is to evaluate or rigorously constrain the signed profile in (H)--(I) \emph{together with} the complementary returns in (J), through their actual Gamma recurrence, conductor division and coefficient maps. An arbitrary matrix model, a dimension fraction, or the known total \(T\) does not calculate this expression. Conversely, the fact that the expression is not yet evaluated is not an obstruction theorem for the entire cohomological programme.

\section{Why an allocation result must pass through the whole action}\label{why-an-allocation-result-must-pass-through-the-whole-action}

Return to the complete arithmetic quotient \(E\), its multiplication \(A=M_S\) and its metric \(G_N\). The monic arithmetic recurrence, in the original coordinate \(S=c+iy\), is

\[
Sp_n=p_{n+1}+cp_n-\frac{\omega_n}{\omega_{n-1}}p_{n-1}.
\]

It gives the rank-two identity

\[
AG_N^{-1}+G_N^{-1}A^*-kG_N^{-1}
=\frac{b_{N+1}b_N^*+b_Nb_{N+1}^*}{\omega_N}.
\]

Reflection \(P(S)\mapsto P(k-S)\) is an isometry of the original source and quotient. It forces the total adjacent pairing to vanish. The relative Hermitian action therefore has only two possibly nonzero eigenvalues \(\epsilon_N,-\epsilon_N\). In an orthogonal kernel/image split, the individual mixed pairings need not vanish: they cancel in their sum, and the mixed exterior norm retains the term \(2|z_K|^2\). {[}R14 CAI6--9, p.623.{]}

The exact same-class exterior argument gives

\[
\frac{\delta q(k+1)}2\le\epsilon_N\qquad(N\ge q-1).
\tag{K}
\]

This is a lower benchmark from the hypothetical quartet itself. It is not an independent contradiction. To contradict that quartet by this route, a proved upper control would have to conflict with the lower benchmark on the same original object and regime.

Put \(d_n=\det G_n/\det G_{n-1}\), so \(0<d_n<1\) at the relevant rows. The adjacent formulas retain both factors:

\[
\epsilon_{q-1}^2=\frac{\omega_q}{\omega_{q-1}}(d_q^{-1}-1),
\] \[
\epsilon_N^2=\frac{\omega_{N+1}}{\omega_N}
           (1-d_N)(d_{N+1}^{-1}-1)\quad(N\ge q).
\]

After taking the two original windows, with the original weights, the action is exactly

\[
J_k^{\rm action}
=\mathcal B_k^{\rm ar}+\mathcal W_k^{\rm ar}
 -\mathcal P_- -\mathcal P_+,
\tag{L}
\] \[
\mathcal W_k^{\rm ar}
=\log\frac{\omega_{2q-1}\omega_{2q}}{\omega_{q-1}\omega_q}.
\]

The two nonnegative penalties are the complete weighted sums of \(-\log(1-d_n)\) specified in CAI13. Formula (L) follows by multiplying the adjacent identities: monic-norm ratios telescope, reciprocal \(d_n\) ratios give the four source volumes, and the remaining contraction factors give precisely those penalties. The original kernel and reference boundary contributions cancel through \(F_K+F_B=\mathcal B\). The arithmetic volume return \(\mathcal B_k^{\rm ar}\) remains in (L); it is not a coordinate constant. {[}R14 CAI11--15, §198 pp.624--625.{]}

This explains what a future result can claim. A calculation of the intrinsic mixed allocation is meaningful geometry. Its effect on the original action must be obtained by substitution into (L), including its cancellation and every remaining term. It cannot be promoted to an action upper bound by dropping the term that cancels it.

\section{The growing-degree result already proved}\label{the-growing-degree-result-already-proved}

The preceding 692-page addition concerns the complete relation polynomial. In the real coordinate define

\[
Q(y)=i^{-q}\chi(c+iy),\qquad R_k=k\sqrt{\delta^2+\gamma^2}.
\]

Its roots occur in opposite pairs, with their full multiplicities, and lie in the disk of radius \(R_k\). For every polynomial \(P\) of degree at most \(r\), compare

\[
H_0(P)=\int_{\mathbb R}|y|^{2q}|P(y)|^2\,d\sigma(y),\qquad
H_Q(P)=\int_{\mathbb R}|Q(y)P(y)|^2\,d\sigma(y).
\]

The root-sensitive full-form theorem gives explicit two-sided comparison, including the complete inner region on both half-lines. In its exact notation

\[
\alpha=\pi/2,\quad s_q=2q+1,\quad
Y_{q,r}=\frac{s_q^2}{32\alpha(r+s_q)},\quad
\rho_{q,r}=R_k/Y_{q,r},
\] \[
\eta_{q,r}=\frac{c_+}{c_-}
 \left(1+\frac{4r+4q+2}{\alpha}\right)s_q(3/8)^{s_q},
\quad c_-=\sqrt2e^{-7/6},\quad c_+=\sqrt{2\sqrt5}e^{2/3}.
\]

Under its actual finite guards \(2R_k\le Y_{q,r}\), \(\eta_{q,r}<1\), it proves

\[
(1-\rho_{q,r}^2)^q(1-\eta_{q,r})H_0(P)
\le H_Q(P)
\le\bigl((1+\rho_{q,r}^2)^q+\eta_{q,r}\bigr)H_0(P).
\tag{M}
\]

These are proved conditions with a nonempty verified arithmetic asymptotic range, not missing hypotheses inserted for the desired conclusion. The proof controls all directions through classical Laguerre moment/evaluation bounds \cite{Szego1975,DLMF18Current}, the paired-root estimate and a direct inner-region calculation. It then takes the original monic minima and the original affine quotient minima. The actual relation sequence remains

\[
0\longrightarrow\mathcal P_r\xrightarrow{\times\chi}
\mathcal P_{q+r}\xrightarrow{J_{q+r}}E\longrightarrow0.
\]

The map \(P(S)\mapsto i^{-r}P(c+iy)\) on monic degree-\(r\) polynomials retains its phase and inverse; the entire unit in the arithmetic cohomology map is unchanged. {[}R14 RGC1--8, RGC18--24, §217 pp.685--690; RGI1--3, §215 p.682.{]}

A separately proved all-degree monic reference bound has signal

\[
F_{q,r}=\frac{q^2}{4(q+r)\log(8(q+r)/q)}.
\]

Combining that bound with (M) and the full arithmetic source comparison yields canonical allowance lower bounds for \(r_k=\lfloor qk^\theta\rfloor\), for each fixed \(0<\theta<2/3\) at \(m=1\), and \(0<\theta<4/3\) at each fixed \(m\ge2\). The further logarithmic refinement sets

\[
(d,A_m^{\deg})=
\begin{cases}(2,1),&m=1,\\(3,m-1),&m\ge2,\end{cases}
\quad\theta_*=2(d-1)/3,
\] \[
r_k=\left\lfloor\frac{qk^{\theta_*}}{(\log k)^\zeta}\right\rfloor,
\qquad\zeta>1/3,
\]

and proves

\[
\log\min_{q-1\le N\le q+r_k-1}\epsilon_N
\ge(1-o(1))\frac{A_m^{\deg}}{8\theta_*}
 k^{(d+2)/3}(\log k)^{\zeta-1}.
\tag{N}
\]

The proof retains the floor, the original \(q=e_k(k+1)^2\), the full root radius, and the exact adjacent factors. The new endpoints are \(q-1,q,q+r_k-1,q+r_k\); the old high endpoints cannot be reused after changing \(r\). The root-error-to-signal calculation is \(O((\log k)^{1-3\zeta})+o(1)\), which explains the strict condition \(\zeta>1/3\). That earlier estimate established neither \(\zeta=1/3\) nor the uncorrected pure-power endpoint; the successor in Section~\ref{sec:moving-cutoff} proves a larger range. {[}R14 RGC16--32, pp.688--692; RGI11--14, §213 pp.673--675; RGI8--10, §198 pp.629--630.{]}

These canonical lower bounds eventually exceed the benchmark (K) by a ratio tending to infinity in their stated ranges. They locate the behavior of this specified action construction. They do not evaluate the independent mixed profile in (H)--(J), and they do not close the proposed tau/Deligne programme. A further analytic extension must prove its error for these same forms and compute its effect through (L).

\section{The accepted moving-cutoff successor}
\label{sec:moving-cutoff}

The accepted 726-page edition extends the previous comparison to degrees
whose ratio to the relation degree grows. This extension is useful because
the original question concerns increasingly large spaces of representatives.
The quotient, the complete zero multiplicities, and the arithmetic norm all
stay fixed as defined above. The degree of a representative is allowed to
increase, and the estimates must control that increase uniformly.
Write \textbf{R726} for this edition; its complete new proofs are
GRF1--27 (pp.~698--705), GRE1--43 (pp.~706--712), GEL1--13
(pp.~713--717), and GRA1--36 (pp.~718--726).

The analytic comparison uses the original Gamma measure and classical
Laguerre and elliptic-integral formulas. Their human sources are
\citet{Szego1975}, the authors of DLMF Chapter~18
\citep{DLMF18Current}, and B.~C.~Carlson's DLMF Chapter~19
\citep{DLMF19Current}. The parameter-dependent estimates are proved in
GRF, GRE and GEL; citing those classical formulas does not attribute
the programme's arithmetic transfer to their authors.

\subsection{The scalar ratio measures a constraint on the original source}
For an integer \(j\ge0\), let \(\omega_d^{\rm ar}\) be the least
squared arithmetic norm of a monic degree-\(d\) polynomial, and put
\[
 \nu_j^{\rm ar}=
 \min_{P\text{ monic},\,\deg P=j}\|\chi P\|_{\rm ar}^2,
 \qquad T_j=\nu_j^{\rm ar}/\omega_{q+j}^{\rm ar},
 \qquad \beta_d=\omega_d^{\rm ar}/\omega_{d-1}^{\rm ar}.
\]
Both minima are on their displayed affine spaces with the original
arithmetic norm. Every \(\chi P\) in the first minimum is monic of
degree \(q+j\), but its required divisibility restricts the available
representatives. The ratio records the cost of that restriction.
The quotient determinant calculation gives
\[
 d_{q+j}=T_j^{-1},\qquad
 \epsilon_{q+j-1}^2=\beta_{q+j}(T_j-1)
 \begin{cases}1,&j=0,\\1-T_{j-1}^{-1},&j\ge1.
 \end{cases}
\tag{O}
\]
To see the first identity, use the source frame
\((\chi,\chi S,\ldots,\chi S^j;1,S,\ldots,S^{q-1})\).
Its determinant has modulus one after arranging its monic columns by
degree. Orthogonal elimination therefore factors the source Gram
determinant into the relation Gram determinant and the quotient
determinant. Taking consecutive ratios yields
\(d_{q+j}=\omega_{q+j}^{\rm ar}/\nu_j^{\rm ar}\).
Substitution into the adjacent formulas of the preceding section proves
the second identity, including the first-row exception. The complete
source calculation is R726 GRA9--17, pp.~719--721.

\subsection{The larger degree range and its exact scale}
Retain \(k=4\ell+1\), \(q=[1+k(m-1)](k+1)^2\), and the fixed
complete quartet. Let \(r_k\) be an integer sequence and set
\[
 t_k=r_k/q\longrightarrow\infty,
 \qquad t_k\log t_k=o(q/k),
 \qquad H_k=\frac{q}{t_k\log t_k}.
\]
The accepted result, including the whole earlier band rather than just
its last row, is
\[
 \min_{0\le j\le r_k}\log T_j\sim H_k,
 \qquad
 \log\min_{q-1\le N\le q+r_k-1}\epsilon_N\sim H_k/2.
\tag{P}
\]
Here \(H_k/k\to\infty\). The moving comparison first bounds the
original full forms on every coefficient vector. Taking their minima
on the same monic and quotient fibres then proves the scalar transfer.
Its moving root error and its arithmetic error are controlled throughout
the specified range. Small degrees, compact positive degree ratios, and
growing ratios are joined by the complete uniform argument in R726
GRB5--8, pp.~678--679; the scalar providers are GRA18--25.
Thus (P) is a theorem of the stated source edition, with an explicit
range, not an inference from a numerical fit.

For a concrete consequence, retain the earlier pair
\((d,A_m^{\deg})=(2,1)\) for \(m=1\) and \((3,m-1)\) for fixed
\(m\ge2\). For
\[
 r_k=\left\lfloor\frac{qk^\theta}{(\log k)^\zeta}\right\rfloor,
\]
one obtains
\[
 \log\min_{q-1\le N\le q+r_k-1}\epsilon_N
 \sim\frac{A_m^{\deg}}{2\theta}
      k^{d-\theta}(\log k)^{\zeta-1}
\tag{Q}
\]
when \(0<\theta<d-1\), for any fixed real \(\zeta\), or when
\(\theta=d-1\) and \(\zeta>1\).
Indeed the floor gives
\(t_k=k^\theta(\log k)^{-\zeta}+O(1/q)\), whence
\[
 \frac{t_k\log t_k}{q/k}
 \sim\frac{\theta}{A_m^{\deg}}
       k^{\theta-d+1}(\log k)^{1-\zeta}.
\]
This tends to zero in exactly the displayed cases, and substitution
into (P) proves (Q). The earlier pure-power ranges \(\theta<2/3\)
and \(\theta<4/3\) have therefore been enlarged to \(\theta<1\)
and \(\theta<2\), respectively, with the stated logarithmic
approach to the new endpoints.

\subsection{A separately proved enclosure at the critical coefficient}
The equality case \(\theta=d-1,\zeta=1\) needs its actual
arithmetic error. Define from the original density
\[
 \vartheta_h=\int_{-1}^1w_h(y)\,dy,\qquad
 M_{\pi/2}=\int_{\mathbb R}e^{\pi y/2}w_h(y)\,dy,
 \qquad\mathfrak s_h=\frac\pi2+
       \log\frac{M_{\pi/2}}{\vartheta_h}.
\]
The proved tilted tail bound makes \(M_{\pi/2}\) finite.
Evenness gives
\[
M_{\pi/2}=\int\cosh(\pi y/2)w_h(y)\,dy
\ge\int w_h\ge\vartheta_h>0,
\]
so \(\mathfrak s_h\ge\pi/2\).
For a fixed coefficient satisfying
\[
 0<c_{\deg}<\frac{A_m^{\deg}}{2(d-1)\mathfrak s_h},\quad
 r_k=\left\lfloor\frac{c_{\deg}qk^{d-1}}{\log k}\right\rfloor,
 \quad B=\frac{A_m^{\deg}}{c_{\deg}(d-1)},
\]
R726 GRB11--13, pp.~679--680, proves
\[
 \frac B2-\mathfrak s_h
 \le\liminf\frac1k\log\min_{q-1\le N\le q+r_k-1}\epsilon_N
 \le\limsup\frac1k\log\min_{q-1\le N\le q+r_k-1}\epsilon_N
 \le\frac B2+\mathfrak s_h.
\tag{R}
\]
All limits use the original subsequence \(k=4\ell+1\).
The lower endpoint is positive. This is an enclosure; the earlier
asymptotic equivalent (P) is not asserted at this boundary. The coefficient
\(c_{\deg}\) is independent of \(k\) and is distinct from the
unchanged polynomial centre \(c=k/2\).

\subsection{The complete shifted action keeps every endpoint}
For an integer shift \(s\ge1\), put \(a=q+s-1\), \(b=a+q\),
and \(w_0=w_q=1\), \(w_i=2\) for \(1\le i<q\). Then
\[
 J_{k,s}^{\rm action}=2\sum_{i=0}^q w_i\log\epsilon_{a+i}
 =\log\frac{V_aV_{a+1}}{V_bV_{b+1}}
  +\log\frac{\omega_b\omega_{b+1}}{\omega_a\omega_{a+1}}
  -\mathcal P_{k,s},
\tag{S}
\]
where \(V_N=\det G_N\), \(\omega_N=\omega_N^{\rm ar}\), and
\[
 \mathcal P_{k,s}=
 \sum_{j=0}^{q+1}c_j[-\log(1-T_{s-1+j}^{-1})],\qquad
 (c_0,\ldots,c_{q+1})=(1,3,\underbrace{4,\ldots,4}_{q-2},3,1).
\]
For any positive sequence \(z_N\), collecting coefficients gives
\[
 \sum_{i=0}^q w_i(\log z_{a+i+1}-\log z_{a+i})
 =\log z_b+\log z_{b+1}-\log z_a-\log z_{a+1}.
\]
Apply this to \(V\) and \(\omega\) in (O). The two adjacent
contractions give \(c_j=w_j+w_{j-1}\), with absent weights zero,
and therefore total penalty weight \(4q\). This proves (S).
R726 GRA27--36 and GRB9--10 supply all finite error bounds and the
transfer through the original kernel inclusion and quotient map.

For \(t=s/q\to\infty\) with \(t\log t=o(q/k)\), those bounds give
\[
 J_{k,s}^{\rm action}\sim\frac{2q^2}{t\log t},\qquad
 \log\mathcal P_{k,s}\sim-\frac{q}{t\log t}.
\]
At the critical choice \(s=r_k\) in (R), the separate enclosure is
\[
 2(B-\mathfrak s_h)\le
 \liminf\frac{J_{k,s}^{\rm action}}{qk}\le
 \limsup\frac{J_{k,s}^{\rm action}}{qk}\le2(B+\mathfrak s_h).
\]
These advances concern the complete correlated action. The intrinsic
signed eight-matrix allocation and each CPQ complementary return still
require their own calculation on the maps already defined above.


\section{A usable reading route and credit map}\label{a-usable-reading-route-and-credit-map}

The following route leads from the original objects to their exact maps, the current estimates and the remaining calculation:

{
\begin{longtable}[]{@{}
  >{\raggedright\arraybackslash}p{\dimexpr.5\linewidth-\tabcolsep\relax}
  >{\raggedright\arraybackslash}p{\dimexpr.5\linewidth-\tabcolsep\relax}@{}}
\toprule\noalign{}
\begin{minipage}[b]{\linewidth}\raggedright
Question
\end{minipage} & \begin{minipage}[b]{\linewidth}\raggedright
Exact proof location in the fixed edition
\end{minipage} \\
\midrule\noalign{}
\endhead
\bottomrule\noalign{}
\endlastfoot
What are the quartet, multiplicities and original tensor map? & R09 §1; R14 OCS1--3, §1 pp.20--21; CAI1 and CAI10, §198 pp.622,624. \\
What is the actual period/cyclic observation? & R14 OCS2--21, §1 pp.20--25; OQX3 §3 pp.32--33; PQO/RPC §§4--5 pp.34--42. \\
How is a representative change proved to preserve the cohomology class? & R14 AKP1--5, §26 pp.110--111; full analytic and categorical foundations in Appendix~\ref{app:foundational-maps}. \\
Where are all full-unit and confluent-jet factors retained? & R09 §2; R14 OCS12--18, §1 pp.23--25; RGI1--3 §215 p.682. \\
How are arithmetic and Gamma metrics related? & R09 §§2--3; R14 AMT §8 pp.52ff.; NSO/EOR §197 pp.602--620; EFP §§204--211 pp.648--659. \\
Where are the actual minimum metrics, frames and four signs proved? & R14 AKP6--10 §26 pp.111--112; CAI2--5 §198 pp.622--623; HJR14--22 §73 pp.304--309. \\
What total is known, and why does it not give the two allocations? & R14 ECL15--18 §59 pp.243--244; HJR16--25 §73 pp.304--309; current prompt, ``Exact working quantities.'' \\
Where is the mixed profile connected to the original complex? & R14 §§151--159 pp.466--487; CPR/BTK/RPI/RRC §§200--203 pp.633--647; CPQ §194 pp.587--591. \\
What exactly remains to estimate? & R14 SFB §195 pp.592--598, NTE1--7 §199 pp.631--632, EFP §§204--211; equations (H)--(J) here. \\
Why can an allocation bound cancel in the action? & R14 CAI11--15 §198 pp.624--625; changed-window RGI8--10 pp.629--630. \\
Where are the preceding growing-degree estimates? & R14 RGC1--32 §217 pp.685--692 and RGI1--14. \\
What is the accepted moving-cutoff advance? & R726 GRF/GRE/GEL/GRA pp.698--726 and GRB1--13 pp.631--632,674--675,678--680; Section~\ref{sec:moving-cutoff} here. \\
\end{longtable}
}
